\chapter{数据布局、缓存与分配}

\begin{introduction}
  \item 用缓存行和局部性解释真实访问成本
  \item 比较 AoS 与 SoA，而不预设结论
  \item 识别分配、对象大小和间接访问
  \item 选择测量后才值得引入的内存优化
\end{introduction}

\section{复杂度相同，耗时仍可不同}

两个 \(O(n)\) 循环可能相差很大。CPU 以缓存行为单位搬运数据，连续访问允许硬件预取；指针追逐、过大的对象和随机访问会增加缓存未命中。大 O 描述规模增长，不描述常数、布局或尾延迟。

在热循环中，先列出真正读取的字段。如果计算只需要价格和数量，但每个元素还携带字符串、时间戳和元数据，Array of Structures（AoS）会把无关字节一起搬入缓存。Structure of Arrays（SoA）把字段分列，适合扫描少数字段；但处理完整单条事件时，AoS 可能更自然且局部。

\lstinputlisting[caption={对同一工作负载比较 AoS 与 SoA}]{project/benchmarks/layout_benchmark.cpp}

这个程序只是一项本机微基准：一次测量、一个固定数据规模、没有 CPU 固定频率或预热控制。它验证两种实现计算相同 checksum，并展示如何获得初步信号；不能据此发布跨机器结论。

\section{对象大小与对齐}

成员顺序会引入填充。使用 \texttt{sizeof} 和 \texttt{alignof} 观察类型，而不要凭字段和相加猜测。把常用小字段集中可能缩小对象，但为了省几个字节破坏清晰语义并不总值得。网络或磁盘协议布局应显式序列化，不能把内存结构直接当稳定格式。

\section{分配是系统行为}

动态分配包含查找可用块、同步、元数据和潜在缺页。\texttt{vector::reserve} 可减少可预测增长中的重复分配。对象池和 \texttt{std::pmr} 资源适合分配模式已知且测得分配热点的场景，但会引入生命周期、峰值容量和线程归属问题。

短生命周期批量对象可以使用 arena，在整批结束时统一释放。前提是任何引用都不能逃出 arena 生命周期；否则“更快释放”换来了悬空风险。

\section{false sharing}

两个线程写不同变量，如果变量落在同一缓存行，缓存一致性仍会在核心间来回转移该行。按线程分片、减少共享写入或适当填充可以缓解。盲目添加 \texttt{alignas(64)} 会增大内存并假设硬件缓存行，应在目标平台测量。

\begin{interviewcheck}
AoS 与 SoA 各适合什么访问模式？为什么链表的常数时间插入不保证更快？对象池增加了哪些正确性风险？什么是 false sharing？
\end{interviewcheck}

\section{练习}

\begin{enumerate}
  \item 多次运行布局基准，报告中位数和离散程度，不只保留最好一次。
  \item 删除 AoS 的时间戳字段后重测，解释对象大小与结果变化。
  \item 用 profiler 验证分配是否真是 CSV 读取热点，再决定是否引入预留容量。
  \item 设计按线程分片的计数器并说明最终归约时机。
\end{enumerate}

\section{本章小结}

性能来自访问模式与硬件层次的互动。连续布局、减少无关字节和控制分配是有力工具，但都要以工作负载和测量为依据。
